\documentclass[12pt]{extarticle}

% ===== Encodage & langue =====
\usepackage{fontspec}
\usepackage{polyglossia}
\setdefaultlanguage{french}

% ===== Police Verdana (compiler avec XeLaTeX ou LuaLaTeX) =====
\setsansfont{Verdana}
\renewcommand{\familydefault}{\sfdefault}

% ===== Interligne =====
\usepackage{setspace}
\setstretch{1.1}

% ===== Marges réduites =====
\usepackage{geometry}
\geometry{a4paper, top=1.8cm, bottom=1.8cm, left=2cm, right=2cm}

% ===== Mathématiques =====
\usepackage{amsmath}

% ===== Espacement réduit autour des équations =====
\setlength{\abovedisplayskip}{4pt}
\setlength{\belowdisplayskip}{4pt}
\setlength{\abovedisplayshortskip}{2pt}
\setlength{\belowdisplayshortskip}{2pt}

% ===== Espacement réduit des sections =====
\usepackage{titlesec}
\titlespacing*{\section}{0pt}{1.2ex}{0.6ex}

% ===== Pas de numéro de page =====
\pagestyle{empty}

\begin{document}

\begin{center}
\Large\textbf{Correction du contrôle séquence 1 version 1}
\end{center}

% --------------------------------------------------
\section*{Exercice 1}
\begin{align*}
18 - 6 - 2 + 19 - 3 &= 12 - 2 + 19 - 3 \\
                     &= 10 + 19 - 3 \\
                     &= 29 - 3 \\
                     &= 26
\end{align*}

% --------------------------------------------------
\section*{Exercice 2}

\begin{align*}
\textbf{a)}
5 + 3 \times 2 &= 5 + 6 \\
               &= 11
\end{align*}

\textbf{b)}
\begin{align*}
34 - 6 \times 2 &= 34 - 12 \\
                &= 22
\end{align*}

\textbf{c)}
\begin{align*}
29 + 12 \div 4 &= 29 + 3 \\
               &= 32
\end{align*}

\textbf{d)}
\begin{align*}
6 + 9 \times 3 + 2 - 4 \div 4 &= 6 + 27 + 2 - 1 \\
                              &= 34
\end{align*}

% --------------------------------------------------
\section*{Exercice 3}
\begin{align*}
3 \times 7 + 2 \times 8 &= 21 + 16 \\
                        &= 37
\end{align*}

La dépense est de 37~€.

% --------------------------------------------------
\section*{Exercice 4}

\textbf{a)} $27 - (6 + 4) \div 2$
\begin{align*}
&= 27 - 10 \div 2 \\
&= 27 - 5 \\
&= 22
\end{align*}

\textbf{b)} $(3 + 8) \times 3 - 7 + 5$
\begin{align*}
&= 11 \times 3 - 7 + 5 \\
&= 33 - 7 + 5 \\
&= 31
\end{align*}

\textbf{c)} $(19 - (5 + 7)) \times 3$
\begin{align*}
&= (19 - 12) \times 3 \\
&= 7 \times 3 \\
&= 21
\end{align*}

\textbf{d)} $4 \times ((21 - (6 + 13)) - 7$
\begin{align*}
&= 4 \times (21 - 19) - 7 \\
&= 4 \times 2 - 7 \\
&= 8 - 7 \\
&= 1
\end{align*}

% --------------------------------------------------
\section*{Exercice 5}
\begin{align*}
11 - 4 \times 2 &= 11 - 8 \\
                &= 3
\end{align*}

Une gomme coûte 3~€.

% --------------------------------------------------
\section*{Exercice bonus 1}

\textbf{a)} $7 \times 3 + 13$

\textbf{b)} $9 \times (7 + 1)$

% --------------------------------------------------
\section*{Exercice bonus 2}

\textbf{a)} La somme du produit de 8 par 3 et de 7.

\textbf{b)} Le produit de la somme de 6 et de 5 par 4.

\newpage

\begin{center}
\Large\textbf{Correction du contrôle séquence 1 version 2}
\end{center}

% --------------------------------------------------
\section*{Exercice 1}
\begin{align*}
17 - 4 - 5 + 23 - 2 &= 13 - 5 + 23 - 2 \\
                     &= 8 + 23 - 2 \\
                     &= 31 - 2 \\
                     &= 29
\end{align*}

% --------------------------------------------------
\section*{Exercice 2}

\begin{align*}
\textbf{a)}
8 + 4 \times 3 &= 8 + 12 \\
               &= 20
\end{align*}

\textbf{b)}
\begin{align*}
41 - 7 \times 3 &= 41 - 21 \\
                &= 20
\end{align*}

\textbf{c)}
\begin{align*}
37 + 15 \div 5 &= 37 + 3 \\
               &= 40
\end{align*}

\textbf{d)}
\begin{align*}
5 + 8 \times 4 + 3 - 6 \div 3 &= 5 + 32 + 3 - 2 \\
                              &= 38
\end{align*}

% --------------------------------------------------
\section*{Exercice 3}
\begin{align*}
4 \times 6 + 2 \times 9 &= 24 + 18 \\
                        &= 42
\end{align*}

La dépense est de 42~€.

% --------------------------------------------------
\section*{Exercice 4}

\textbf{a)} $32 - (5 + 9) \div 2$
\begin{align*}
&= 32 - 14 \div 2 \\
&= 32 - 7 \\
&= 25
\end{align*}

\textbf{b)} $(4 + 7) \times 4 - 9 + 3$
\begin{align*}
&= 11 \times 4 - 9 + 3 \\
&= 44 - 9 + 3 \\
&= 38
\end{align*}

\textbf{c)} $(21 - (8 + 5)) \times 4$
\begin{align*}
&= (21 - 13) \times 4 \\
&= 8 \times 4 \\
&= 32
\end{align*}

\textbf{d)} $5 \times (25 - (9 + 11)) - 7$
\begin{align*}
&= 5 \times (25 - 20) - 7 \\
&= 5 \times 5 - 7 \\
&= 25 - 7 \\
&= 18
\end{align*}

% --------------------------------------------------
\section*{Exercice 5}
\begin{align*}
14 - 5 \times 2 &= 14 - 10 \\
                &= 4
\end{align*}

Une gomme coûte 4~€.

% --------------------------------------------------
\section*{Exercice bonus 1}

\textbf{a)} $8 \times 4 + 11$

\textbf{b)} $7 \times (9 + 2)$

% --------------------------------------------------
\section*{Exercice bonus 2}

\textbf{a)} La somme du produit de 9 par 4 et de 6.

\textbf{b)} Le produit de la somme de 7 et de 4 par 5.

\newpage

\begin{center}
\Large\textbf{Correction du contrôle séquence 1 version 3}
\end{center}

% --------------------------------------------------
\section*{Exercice 1}
\begin{align*}
15 - 3 - 4 + 25 - 5 &= 12 - 4 + 25 - 5 \\
                     &= 8 + 25 - 5 \\
                     &= 33 - 5 \\
                     &= 28
\end{align*}

% --------------------------------------------------
\section*{Exercice 2}

\begin{align*}
\textbf{a)}
7 + 5 \times 4 &= 7 + 20 \\
               &= 27
\end{align*}

\textbf{b)}
\begin{align*}
46 - 8 \times 4 &= 46 - 32 \\
                &= 14
\end{align*}

\textbf{c)}
\begin{align*}
43 + 18 \div 6 &= 43 + 3 \\
               &= 46
\end{align*}

\textbf{d)}
\begin{align*}
8 + 7 \times 5 + 4 - 8 \div 2 &= 8 + 35 + 4 - 4 \\
                              &= 43
\end{align*}

% --------------------------------------------------
\section*{Exercice 3}
\begin{align*}
3 \times 8 + 3 \times 7 &= 24 + 21 \\
                        &= 45
\end{align*}

La dépense est de 45~€.

% --------------------------------------------------
\section*{Exercice 4}

\textbf{a)} $36 - (7 + 9) \div 2$
\begin{align*}
&= 36 - 16 \div 2 \\
&= 36 - 8 \\
&= 28
\end{align*}

\textbf{b)} $(5 + 9) \times 4 - 8 + 6$
\begin{align*}
&= 14 \times 4 - 8 + 6 \\
&= 56 - 8 + 6 \\
&= 54
\end{align*}

\textbf{c)} $(23 - (6 + 8)) \times 5$
\begin{align*}
&= (23 - 14) \times 5 \\
&= 9 \times 5 \\
&= 45
\end{align*}

\textbf{d)} $6 \times (27 - (8 + 14)) - 5$
\begin{align*}
&= 6 \times (27 - 22) - 5 \\
&= 6 \times 5 - 5 \\
&= 30 - 5 \\
&= 25
\end{align*}

% --------------------------------------------------
\section*{Exercice 5}
\begin{align*}
16 - 6 \times 2 &= 16 - 12 \\
                &= 4
\end{align*}

Une gomme coûte 4~€.

% --------------------------------------------------
\section*{Exercice bonus 1}

\textbf{a)} $9 \times 5 + 12$

\textbf{b)} $6 \times (8 + 3)$

% --------------------------------------------------
\section*{Exercice bonus 2}

\textbf{a)} La somme du produit de 7 par 5 et de 8.

\textbf{b)} Le produit de la somme de 8 et de 3 par 6.

\newpage

\begin{center}
\Large\textbf{Correction du contrôle séquence 1 version 4}
\end{center}

% --------------------------------------------------
\section*{Exercice 1}
\begin{align*}
19 - 5 - 3 + 21 - 6 &= 14 - 3 + 21 - 6 \\
                     &= 11 + 21 - 6 \\
                     &= 32 - 6 \\
                     &= 26
\end{align*}

% --------------------------------------------------
\section*{Exercice 2}

\begin{align*}
\textbf{a)}
6 + 4 \times 5 &= 6 + 20 \\
               &= 26
\end{align*}

\textbf{b)}
\begin{align*}
38 - 5 \times 4 &= 38 - 20 \\
                &= 18
\end{align*}

\textbf{c)}
\begin{align*}
31 + 16 \div 8 &= 31 + 2 \\
               &= 33
\end{align*}

\textbf{d)}
\begin{align*}
7 + 6 \times 4 + 5 - 9 \div 3 &= 7 + 24 + 5 - 3 \\
                              &= 33
\end{align*}

% --------------------------------------------------
\section*{Exercice 3}
\begin{align*}
4 \times 6 + 3 \times 7 &= 24 + 21 \\
                        &= 45
\end{align*}

La dépense est de 45~€.

% --------------------------------------------------
\section*{Exercice 4}

\textbf{a)} $29 - (4 + 8) \div 2$
\begin{align*}
&= 29 - 12 \div 2 \\
&= 29 - 6 \\
&= 23
\end{align*}

\textbf{b)} $(6 + 7) \times 5 - 8 + 4$
\begin{align*}
&= 13 \times 5 - 8 + 4 \\
&= 65 - 8 + 4 \\
&= 61
\end{align*}

\textbf{c)} $(25 - (9 + 6)) \times 4$
\begin{align*}
&= (25 - 15) \times 4 \\
&= 10 \times 4 \\
&= 40
\end{align*}

\textbf{d)} $5 \times (24 - (7 + 12)) - 6$
\begin{align*}
&= 5 \times (24 - 19) - 6 \\
&= 5 \times 5 - 6 \\
&= 25 - 6 \\
&= 19
\end{align*}

% --------------------------------------------------
\section*{Exercice 5}
\begin{align*}
13 - 4 \times 3 &= 13 - 12 \\
                &= 1
\end{align*}

Une gomme coûte 1~€.

% --------------------------------------------------
\section*{Exercice bonus 1}

\textbf{a)} $6 \times 5 + 14$

\textbf{b)} $8 \times (6 + 3)$

% --------------------------------------------------
\section*{Exercice bonus 2}

\textbf{a)} La somme du produit de 6 par 5 et de 9.

\textbf{b)} Le produit de la somme de 9 et de 2 par 7.

\newpage

\begin{center}
\Large\textbf{Correction du contrôle séquence 1 version 5}
\end{center}

% --------------------------------------------------
\section*{Exercice 1}
\begin{align*}
16 - 7 - 1 + 24 - 4 &= 9 - 1 + 24 - 4 \\
                     &= 8 + 24 - 4 \\
                     &= 32 - 4 \\
                     &= 28
\end{align*}

% --------------------------------------------------
\section*{Exercice 2}

\begin{align*}
\textbf{a)}
9 + 4 \times 5 &= 9 + 20 \\
               &= 29
\end{align*}

\textbf{b)}
\begin{align*}
47 - 9 \times 4 &= 47 - 36 \\
                &= 11
\end{align*}

\textbf{c)}
\begin{align*}
35 + 14 \div 7 &= 35 + 2 \\
               &= 37
\end{align*}

\textbf{d)}
\begin{align*}
9 + 7 \times 4 + 6 - 8 \div 2 &= 9 + 28 + 6 - 4 \\
                              &= 39
\end{align*}

% --------------------------------------------------
\section*{Exercice 3}
\begin{align*}
5 \times 6 + 2 \times 9 &= 30 + 18 \\
                        &= 48
\end{align*}

La dépense est de 48~€.

% --------------------------------------------------
\section*{Exercice 4}

\textbf{a)} $34 - (8 + 6) \div 2$
\begin{align*}
&= 34 - 14 \div 2 \\
&= 34 - 7 \\
&= 27
\end{align*}

\textbf{b)} $(7 + 6) \times 5 - 9 + 4$
\begin{align*}
&= 13 \times 5 - 9 + 4 \\
&= 65 - 9 + 4 \\
&= 60
\end{align*}

\textbf{c)} $(27 - (7 + 9)) \times 4$
\begin{align*}
&= (27 - 16) \times 4 \\
&= 11 \times 4 \\
&= 44
\end{align*}

\textbf{d)} $6 \times (26 - (9 + 11)) - 7$
\begin{align*}
&= 6 \times (26 - 20) - 7 \\
&= 6 \times 6 - 7 \\
&= 36 - 7 \\
&= 29
\end{align*}

% --------------------------------------------------
\section*{Exercice 5}
\begin{align*}
17 - 5 \times 3 &= 17 - 15 \\
                &= 2
\end{align*}

Une gomme coûte 2~€.

% --------------------------------------------------
\section*{Exercice bonus 1}

\textbf{a)} $7 \times 5 + 15$

\textbf{b)} $6 \times (9 + 4)$

% --------------------------------------------------
\section*{Exercice bonus 2}

\textbf{a)} La somme du produit de 7 par 6 et de 5.

\textbf{b)} Le produit de la somme de 5 et de 8 par 6.

\newpage

\begin{center}
\Large\textbf{Correction du contrôle séquence 1 version 6}
\end{center}

% --------------------------------------------------
\section*{Exercice 1}
\begin{align*}
14 - 2 - 6 + 27 - 5 &= 12 - 6 + 27 - 5 \\
                     &= 6 + 27 - 5 \\
                     &= 33 - 5 \\
                     &= 28
\end{align*}

% --------------------------------------------------
\section*{Exercice 2}

\begin{align*}
\textbf{a)}
4 + 5 \times 6 &= 4 + 30 \\
               &= 34
\end{align*}

\textbf{b)}
\begin{align*}
52 - 7 \times 5 &= 52 - 35 \\
                &= 17
\end{align*}

\textbf{c)}
\begin{align*}
39 + 21 \div 7 &= 39 + 3 \\
               &= 42
\end{align*}

\textbf{d)}
\begin{align*}
4 + 8 \times 5 + 7 - 10 \div 2 &= 4 + 40 + 7 - 5 \\
                              &= 46
\end{align*}

% --------------------------------------------------
\section*{Exercice 3}
\begin{align*}
3 \times 9 + 2 \times 10 &= 27 + 20 \\
                        &= 47
\end{align*}

La dépense est de 47~€.

% --------------------------------------------------
\section*{Exercice 4}

\textbf{a)} $38 - (9 + 7) \div 2$
\begin{align*}
&= 38 - 16 \div 2 \\
&= 38 - 8 \\
&= 30
\end{align*}

\textbf{b)} $(8 + 5) \times 6 - 9 + 7$
\begin{align*}
&= 13 \times 6 - 9 + 7 \\
&= 78 - 9 + 7 \\
&= 76
\end{align*}

\textbf{c)} $(29 - (8 + 9)) \times 6$
\begin{align*}
&= (29 - 17) \times 6 \\
&= 12 \times 6 \\
&= 72
\end{align*}

\textbf{d)} $7 \times (28 - (9 + 13)) - 8$
\begin{align*}
&= 7 \times (28 - 22) - 8 \\
&= 7 \times 6 - 8 \\
&= 42 - 8 \\
&= 34
\end{align*}

% --------------------------------------------------
\section*{Exercice 5}
\begin{align*}
20 - 6 \times 3 &= 20 - 18 \\
                &= 2
\end{align*}

Une gomme coûte 2~€.

% --------------------------------------------------
\section*{Exercice bonus 1}

\textbf{a)} $8 \times 6 + 13$

\textbf{b)} $7 \times (9 + 4)$

% --------------------------------------------------
\section*{Exercice bonus 2}

\textbf{a)} La somme du produit de 8 par 7 et de 6.

\textbf{b)} Le produit de la somme de 4 et de 9 par 7.

\newpage

\begin{center}
\Large\textbf{Correction du contrôle séquence 1 version 7}
\end{center}

% --------------------------------------------------
\section*{Exercice 1}
\begin{align*}
13 - 6 - 3 + 28 - 4 &= 7 - 3 + 28 - 4 \\
                     &= 4 + 28 - 4 \\
                     &= 32 - 4 \\
                     &= 28
\end{align*}

% --------------------------------------------------
\section*{Exercice 2}

\begin{align*}
\textbf{a)}
5 + 6 \times 4 &= 5 + 24 \\
               &= 29
\end{align*}

\textbf{b)}
\begin{align*}
43 - 8 \times 5 &= 43 - 40 \\
                &= 3
\end{align*}

\textbf{c)}
\begin{align*}
45 + 24 \div 8 &= 45 + 3 \\
               &= 48
\end{align*}

\textbf{d)}
\begin{align*}
5 + 7 \times 6 + 8 - 6 \div 2 &= 5 + 42 + 8 - 3 \\
                              &= 52
\end{align*}

% --------------------------------------------------
\section*{Exercice 3}
\begin{align*}
4 \times 7 + 3 \times 8 &= 28 + 24 \\
                        &= 52
\end{align*}

La dépense est de 52~€.

% --------------------------------------------------
\section*{Exercice 4}

\textbf{a)} $33 - (6 + 8) \div 2$
\begin{align*}
&= 33 - 14 \div 2 \\
&= 33 - 7 \\
&= 26
\end{align*}

\textbf{b)} $(4 + 9) \times 5 - 6 + 8$
\begin{align*}
&= 13 \times 5 - 6 + 8 \\
&= 65 - 6 + 8 \\
&= 67
\end{align*}

\textbf{c)} $(26 - (7 + 8)) \times 5$
\begin{align*}
&= (26 - 15) \times 5 \\
&= 11 \times 5 \\
&= 55
\end{align*}

\textbf{d)} $5 \times (23 - (6 + 11)) - 8$
\begin{align*}
&= 5 \times (23 - 17) - 8 \\
&= 5 \times 6 - 8 \\
&= 30 - 8 \\
&= 22
\end{align*}

% --------------------------------------------------
\section*{Exercice 5}
\begin{align*}
14 - 4 \times 3 &= 14 - 12 \\
                &= 2
\end{align*}

Une gomme coûte 2~€.

% --------------------------------------------------
\section*{Exercice bonus 1}

\textbf{a)} $9 \times 4 + 16$

\textbf{b)} $8 \times (7 + 5)$

% --------------------------------------------------
\section*{Exercice bonus 2}

\textbf{a)} La somme du produit de 9 par 5 et de 7.

\textbf{b)} Le produit de la somme de 6 et de 9 par 8.

\newpage

\begin{center}
\Large\textbf{Correction du contrôle séquence 1 version 8}
\end{center}

% --------------------------------------------------
\section*{Exercice 1}
\begin{align*}
12 - 4 - 7 + 29 - 3 &= 8 - 7 + 29 - 3 \\
                     &= 1 + 29 - 3 \\
                     &= 30 - 3 \\
                     &= 27
\end{align*}

% --------------------------------------------------
\section*{Exercice 2}

\begin{align*}
\textbf{a)}
6 + 5 \times 7 &= 6 + 35 \\
               &= 41
\end{align*}

\textbf{b)}
\begin{align*}
56 - 9 \times 5 &= 56 - 45 \\
                &= 11
\end{align*}

\textbf{c)}
\begin{align*}
41 + 27 \div 9 &= 41 + 3 \\
               &= 44
\end{align*}

\textbf{d)}
\begin{align*}
6 + 8 \times 6 + 9 - 8 \div 2 &= 6 + 48 + 9 - 4 \\
                              &= 59
\end{align*}

% --------------------------------------------------
\section*{Exercice 3}
\begin{align*}
3 \times 8 + 4 \times 7 &= 24 + 28 \\
                        &= 52
\end{align*}

La dépense est de 52~€.

% --------------------------------------------------
\section*{Exercice 4}

\textbf{a)} $37 - (5 + 9) \div 2$
\begin{align*}
&= 37 - 14 \div 2 \\
&= 37 - 7 \\
&= 30
\end{align*}

\textbf{b)} $(5 + 8) \times 6 - 7 + 9$
\begin{align*}
&= 13 \times 6 - 7 + 9 \\
&= 78 - 7 + 9 \\
&= 80
\end{align*}

\textbf{c)} $(28 - (9 + 7)) \times 5$
\begin{align*}
&= (28 - 16) \times 5 \\
&= 12 \times 5 \\
&= 60
\end{align*}

\textbf{d)} $6 \times (29 - (8 + 15)) - 7$
\begin{align*}
&= 6 \times (29 - 23) - 7 \\
&= 6 \times 6 - 7 \\
&= 36 - 7 \\
&= 29
\end{align*}

% --------------------------------------------------
\section*{Exercice 5}
\begin{align*}
19 - 5 \times 3 &= 19 - 15 \\
                &= 4
\end{align*}

Une gomme coûte 4~€.

% --------------------------------------------------
\section*{Exercice bonus 1}

\textbf{a)} $7 \times 6 + 17$

\textbf{b)} $9 \times (8 + 5)$

% --------------------------------------------------
\section*{Exercice bonus 2}

\textbf{a)} La somme du produit de 8 par 6 et de 9.

\textbf{b)} Le produit de la somme de 7 et de 6 par 8.

\end{document}
